\documentclass[11pt]{article}
\usepackage[T1]{fontenc}
\usepackage{lmodern}
\usepackage{amsmath,amssymb}
\usepackage[margin=1in]{geometry}
\usepackage[colorlinks=true,citecolor=blue,urlcolor=blue]{hyperref}
\usepackage{microtype}

\title{The quantitative Banach--Saks equality question\\
\large Exact negative answer in later literature}
\author{Literature-status note for \texttt{fa\_banach\_001}}
\date{21 July 2026}

\begin{document}
\maketitle

\begin{center}
\fbox{\parbox{0.9\textwidth}{\textbf{Status: literature already answered.}
Silber's Example 6.2 explicitly answers the final question of
arXiv:1406.0684 negatively.  This is not an original result of the run.}}
\end{center}

\section*{Original question}

Bendov\'a, Kalenda, and Spurn\'y ask in the final question of Section 6,
page 17 of arXiv:1406.0684 \cite{BKS}:
\begin{quote}
Let $X$ be a Banach space and $A\subset X$ a bounded set. Is it necessarily
true that
\[
 \operatorname{wbs}(A)=2\operatorname{sm}(A)=2\operatorname{wus}(A)?
\]
\end{quote}

\section*{Exact later answer}

Silber's arXiv:2111.12773 \cite{Silber}, Example 6.2 on pages 27--28,
equips a Schreier space with the equivalent norm
\[
 \|x\|_* = \max\{\|x^+\|,\|x^-\|\}
\]
and takes $A=\{e_n:n\in\mathbb N\}$.  The example proves
\[
 \operatorname{sm}_1(A)=\frac12,\qquad
 \operatorname{wus}_1(A)=1,\qquad
 \operatorname{wbs}_0(A)=1.
\]
The paragraph immediately following the example records the notation bridge
\[
 \operatorname{wbs}=\operatorname{wbs}_0,
 \qquad \operatorname{sm}=\operatorname{sm}_1,
 \qquad \operatorname{wus}=\operatorname{wus}_1,
\]
quotes the original equality, and explicitly states that Example 6.2 answers
the question negatively.  Therefore
\[
 \operatorname{wbs}(A)=2\operatorname{sm}(A)=1
 \quad\text{but}\quad
 2\operatorname{wus}(A)=2.
\]

This is an exact, complete, author-identified negative answer rather than an
agent-inferred implication.  A neighboring higher-order equality remains open
in the supporting paper, but it is distinct from the 2015 question.

\begin{thebibliography}{9}
\bibitem{BKS}
H. Bendov\'a, O. F. K. Kalenda, and J. Spurn\'y,
\emph{Quantification of the Banach--Saks property},
J. Funct. Anal. 268 (2015), 1733--1754; arXiv:1406.0684.

\bibitem{Silber}
Z. Silber,
\emph{Quantification of Banach--Saks properties of higher orders},
arXiv:2111.12773 (2021), Example 6.2.
\end{thebibliography}

\end{document}
